\documentclass[11pt, a4paper]{article}

\usepackage[T1]{fontenc}
\usepackage[utf8]{inputenc}
\usepackage{tgpagella}
\usepackage[spanish, es-noshorthands]{babel}
\usepackage{amsmath, amssymb, bm}
\usepackage{booktabs}
\usepackage[table, xcdraw]{xcolor}
\usepackage[most]{tcolorbox}
\usepackage[american]{circuitikz}
\usepackage[margin=2.5cm]{geometry}
\usepackage{fancyhdr}

\pagestyle{fancy}
\fancyhf{}
\setlength{\headheight}{16pt}
\lhead{\textbf{UCB Sede Tarija} | Ing. Mecatrónica}
\chead{}
\rhead{IMT-121 Circuitos Electrónicos I | Ejercicios S06}
\lfoot{Prof: M.Sc. Ing. Bernardo Quiroga Turdera}
\rfoot{Página \thepage}
\renewcommand{\headrulewidth}{0.4pt}

\definecolor{ucbblue}{RGB}{0, 51, 102}

\begin{document}

\begin{tcolorbox}[colback=ucbblue!5!white, colframe=ucbblue, title=\textbf{Hoja de Ejercicios Guiada: Polarización DC del BJT}]
    \centering \large \textbf{Análisis de Redes de Polarización Fija}
\end{tcolorbox}

\section*{Ejercicio 1: Reconocimiento de Región}
Clasifica y justifica la región de operación (Corte, Activa Directa o Saturación) para los siguientes casos:
\begin{enumerate}
    \item Unión BE no polarizada en directa, Unión BC no polarizada en directa.
    \item Unión BE polarizada en directa, Unión BC polarizada en inversa.
    \item Unión BE polarizada en directa, Unión BC polarizada en directa.
\end{enumerate}

\section*{Ejercicio 2: Análisis DC Guiado}
Considera el Esquema de Ejercicio 1 (asumir $\beta=100$, $V_{BE}=0.7\,\text{V}$ en conducción):

\begin{center}
\begin{circuitikz}[scale=1, transform shape, thick]
    \draw (0,0) node[npn](Q){};
    \draw (Q.E) node[ground]{};
    % PROTECCIÓN CON LLAVES EN LAS ETIQUETAS: l={...}
    \draw (Q.C) to[R, l_={$R_C = 2.2\,\text{k}\Omega$}] (0,3);
    \draw (Q.B) -- (-2.5, 0) to[R, l={$R_B = 560\,\text{k}\Omega$}] (-2.5, 3);
    \draw (-2.5,3) -- (0,3) -- (0,3.5) node[above]{$+12\,\text{V}$};
\end{circuitikz}
\end{center}

\begin{enumerate}
    \item Añade al esquema (a lápiz) las direcciones de $I_B$, $I_C$, $I_E$, y las polaridades de $V_{BE}$ y $V_{CE}$.
    \item Calcula la corriente de base $I_B$.
    \item Calcula la corriente candidata de colector $I_C$.
    \item Calcula $V_{CE}$.
    \item Comprueba si la suposición de región activa es autoconsistente.
    \item Deriva la ecuación de la recta de carga.
    \item Encuentra las dos intersecciones con los ejes.
    \item Indica el Punto $Q$ exacto.
    \item \textbf{Predicción:} Sin recalcular, ¿qué sucede con $Q$ si $R_B$ aumenta? ¿Qué sucede si $R_C$ aumenta?
\end{enumerate}

\section*{Ejercicio 3: Variante Semi-Guiada}
Utiliza la misma topología de polarización fija del Ejercicio 2, pero en este caso el diagrama no tiene etiquetas de tensión ni corrientes, sólo los componentes en crudo:

\begin{center}
\begin{circuitikz}[scale=1, transform shape, thick]
    \draw (0,0) node[npn](Q){};
    \draw (Q.E) node[ground]{};
    \draw (Q.C) to[R] (0,3);
    \draw (Q.B) -- (-2.5, 0) to[R] (-2.5, 3);
    \draw (-2.5,3) -- (0,3) -- (0,3.5);
\end{circuitikz}
\end{center}

\textbf{Tareas:}
\begin{enumerate}
    \item Dibuja todas las flechas de corriente ($I_B, I_C, I_E$).
    \item Etiqueta explícitamente $V_{BE}$ y $V_{CE}$.
    \item Identifica visualmente la malla base-emisor y la malla colector-emisor.
    \item Escribe las ecuaciones de la Ley de Tensiones de Kirchhoff (KVL) para ambas mallas a partir de este diagrama.
\end{enumerate}

\section*{Ejercicio 4: Problema Independiente}
Analiza el siguiente circuito (Esquema de Ejercicio 2). Datos adicionales: $\beta = 120$, $V_{BE} = 0.7\,\text{V}$.

\begin{center}
\begin{circuitikz}[scale=1, transform shape, thick]
    \draw (0,0) node[npn](Q){};
    \draw (Q.E) node[ground]{};
    % PROTECCIÓN CON LLAVES EN LAS ETIQUETAS: l={...}
    \draw (Q.C) to[R, l_={$R_C = 1.5\,\text{k}\Omega$}] (0,3);
    \draw (Q.B) -- (-2.5, 0) to[R, l={$R_B = 330\,\text{k}\Omega$}] (-2.5, 3);
    \draw (-2.5,3) -- (0,3) -- (0,3.5) node[above]{$+9\,\text{V}$};
\end{circuitikz}
\end{center}

\begin{enumerate}
    \item Calcula $I_B$, $I_C$ y $V_{CE}$.
    \item Decide si la suposición de región activa es válida.
    \item Deriva la recta de carga y encuentra las intersecciones.
    \item Localiza el Punto $Q$.
    \item Predice cualitativamente qué sucede con el circuito si $\beta$ cambia de $120$ a $60$ (por ejemplo, debido a un reemplazo de componente).
\end{enumerate}

\section*{Ejercicio 5: Transferencia Conceptual}
\begin{enumerate}
    \item ¿Por qué $I_C = \beta I_B$ no es una ley universal del transistor?
    \item ¿Cuál es la diferencia entre una curva característica de salida del transistor y una recta de carga del colector?
    \item ¿Qué representa físicamente el punto $Q$?
    \item ¿Por qué es prudente incluir el esquema del circuito explícitamente en los primeros ejercicios de esta sesión en lugar de describirlo únicamente en texto?
\end{enumerate}

\end{document}
